\begin{proof}
For each $n \in \mathbb{N}$, $f_n$ is well-defined. By the Recursion Theorem 6.2 of the Introduction, there exists a unique function 
\[
\varphi : \mathbb{N} \to G
\] 
such that 
\[
\varphi(0) = a_1, \quad \varphi(n+1) = f_n(\varphi(n)) \quad \text{for all } n \in \mathbb{N}.
\]

In particular, 
\[
\varphi(0) = a_1, \quad \varphi(1) = a_1 a_2, \quad \varphi(2) = (a_1 a_2) a_3,
\] 
and for $n \ge 0$, 
\[
\varphi(n+1) = (\varphi(n)) a_{n+1}.
\]

Finally, define 
\[
\psi = \varphi \circ \theta
\] 
with $\theta(k) = k-1$. Then $\psi$ satisfies 
\[
\psi(1) = a_1, \quad \psi(2) = a_1 a_2, \quad \psi(3) = (a_1 a_2) a_3,
\] 
and for $n \ge 1$, 
\[
\psi(n+1) = (\psi(n)) a_{n+1}.
\] 
Uniqueness follows from the Recursion Theorem. 
\end{proof}
